% ============================================================
% Chapter 3: Methodology
% CT-MUSIQ: Automated Perceptual Assessment of Low-Dose CT
% Author: M. Samiul Hasnat, Sichuan University
% ============================================================

\chapter{Methodology}

\section{Overview and Problem Formulation}

The central objective of this thesis is to develop an automated, no-reference image quality
assessment (NR-IQA) system for brain low-dose computed tomography (LDCT) images.
In clinical practice, radiologists manually evaluate each CT scan for diagnostic
acceptability, a task that is time-consuming, subject to inter-rater variability, and
unsuitable for high-throughput quality control pipelines. The proposed system, named
CT-MUSIQ, aims to reproduce this human perceptual judgment automatically by learning
a function that maps an input CT image directly to a continuous quality score, without
requiring a paired reference scan of known perfect quality.

Formally, the task is defined as follows. Given a brain CT slice
$\mathbf{x} \in \mathbb{R}^{H \times W}$ (where $H = W = 512$), the model must predict
a scalar quality score $\hat{q} \in [0, 4]$ that closely matches the consensus radiologist
score $q \in [0, 4]$ assigned on a Likert scale. The model has no access to the original,
noise-free reference image during inference, making this a strictly no-reference problem.
This is clinically important because in routine CT acquisition there is no ``pristine''
reference available---the task is to assess perceptual quality from the single acquired
image alone.

CT-MUSIQ is an architectural adaptation of MUSIQ (Multi-Scale Image Quality
Transformer), a model proposed by Ke et al.\ \cite{ke2021musiq} (ICCV 2021) for natural
image quality assessment. The key innovations introduced in CT-MUSIQ over the vanilla
MUSIQ architecture are: (1) a custom hash-based positional encoding scheme tailored to
multi-scale variable-length patch sequences, (2) per-scale prediction heads that enable
a novel scale-consistency KL divergence regularisation loss, (3) a pairwise margin
ranking loss for improved ordinal discrimination, and (4) a clinically-motivated brain
CT windowing pre-processing step in place of the original challenge's abdominal windowing
parameters. The complete system is trained end-to-end on the LDCTIQAC 2023 challenge
dataset \cite{lee2025ldctiqac} using a two-stage fine-tuning protocol with
mixed-precision training, designed to fit within 6~GB of GPU video memory (VRAM) on a
consumer-grade NVIDIA RTX~3060 Laptop GPU.

The following sections describe each methodological component in detail: the dataset and
pre-processing pipeline (Section~\ref{sec:dataset}), the CT-MUSIQ model architecture
(Section~\ref{sec:architecture}), the composite loss function
(Section~\ref{sec:loss}), the training procedure (Section~\ref{sec:training}), data
augmentation (Section~\ref{sec:augmentation}), the evaluation protocol
(Section~\ref{sec:evaluation}), the comparison baseline (Section~\ref{sec:baseline}),
and the ablation study design (Section~\ref{sec:ablation}). Implementation details and
hardware specifications are provided in Section~\ref{sec:implementation}.


% ============================================================
\section{Dataset and Pre-Processing}
\label{sec:dataset}

\subsection{LDCTIQAC 2023 Dataset}

All experiments in this thesis are conducted on the \textbf{LDCTIQAC 2023 dataset}
(Low-Dose CT Image Quality Assessment Challenge 2023), which was the basis of a
challenge organised at MICCAI 2023 and subsequently documented in Lee et al.\
\cite{lee2025ldctiqac} in \textit{Medical Image Analysis}. The dataset consists of
\textbf{1,000 axial brain CT slices}, each stored as a 32-bit floating-point TIFF image
of dimensions $512 \times 512$ pixels. The pixel values in the distributed files are
pre-normalised to the range $[0, 1]$ prior to release. Each image corresponds to a
single axial slice from a low-dose CT acquisition protocol designed to simulate
clinically realistic noise levels encountered in routine brain CT examinations.

Quality labels were obtained through expert radiologist assessment. Multiple radiologists
independently assigned a perceptual quality score to each image using a five-point
Likert scale ranging from~0 (completely unacceptable, severe noise or artefacts) to~4
(excellent quality, fully diagnostic). The labels distributed with the dataset represent
the averaged consensus score across all annotating radiologists, yielding a continuous
floating-point value for each image. This averaged score is used as the regression
target throughout training and evaluation.

The image files follow a strict four-digit zero-padded naming convention:
\texttt{\{index:04d\}.tif}, yielding filenames \texttt{0000.tif} through
\texttt{0999.tif}. The mapping from filenames to quality scores is provided in a JSON
file (\texttt{train.json}) with the filename (including extension) as the key and the
quality score as the value.

\subsection{Data Splits}

The 1,000 images are partitioned into three non-overlapping subsets based on filename
index order, as specified in the challenge: a \textbf{training set} comprising images
0000 through 0699 (700 images), a \textbf{validation set} comprising images 0700 through
0899 (200 images), and a \textbf{test set} comprising images 0900 through 0999
(100 images). This corresponds to a 70/20/10 split ratio.

Importantly, the split is deterministic and index-based rather than randomly sampled.
This design choice ensures that all reported results are reproducible without knowledge
of a random seed for splitting, and it matches the official challenge protocol. During
training, only the training split is used for gradient updates. The validation split is
used exclusively for hyperparameter tuning, model selection (checkpoint saving based on
best validation aggregate score), and early stopping. The test split is held out
completely until final evaluation and is never used to inform any training decision.

\subsection{CT Windowing Pre-Processing}

CT images store raw pixel values in Hounsfield Units (HU), a quantitative scale where
water is 0~HU, air is approximately $-1000$~HU, and dense bone is approximately
$+1000$~HU. Raw HU values span a large dynamic range (typically $-1024$ to $+3071$~HU)
that is clinically irrelevant for any single tissue type. Radiologists routinely apply a
\textbf{CT window} transformation before viewing: they select a window level (centre)
and window width that clips and linearly rescales the HU range to display only the tissue
of interest with full contrast.

A critical methodological adaptation in this thesis concerns the windowing parameters.
The original LDCTIQAC 2023 challenge was designed around \textbf{abdominal CT} imaging,
for which the clinical standard window is centred at level 40~HU with a width of
350~HU. However, the specific dataset used in this thesis contains \textbf{brain CT
slices}, for which the standard clinical window is the \textbf{brain soft-tissue window}:
level 40~HU, width 80~HU. This choice reflects the practice of neuroradiology: the
narrow brain window (spanning from 0 to 80~HU) optimally differentiates grey matter
($\sim$40~HU), white matter ($\sim$35~HU), cerebrospinal fluid ($\sim$0--15~HU), and
hyperacute haemorrhage ($\sim$50--80~HU), which are the diagnostically relevant tissues
in brain CT quality assessment.

The window transformation is applied as follows. Given an input pixel value $p$ and
window parameters level $L = 40$~HU and width $W = 80$~HU:

\begin{align}
    HU_{\min} &= L - \frac{W}{2} = 40 - 40 = 0~\text{HU} \\
    HU_{\max} &= L + \frac{W}{2} = 40 + 40 = 80~\text{HU} \\
    p_{\text{clipped}} &= \text{clip}(p,\; HU_{\min},\; HU_{\max}) \\
    p_{\text{norm}} &= \frac{p_{\text{clipped}} - HU_{\min}}{HU_{\max} - HU_{\min}} \in [0, 1]
\end{align}

Since the dataset images are already stored as pre-normalised float32 values in $[0,1]$,
the windowing is implicitly applied during dataset generation. No further normalisation
using ImageNet mean and standard deviation is applied, as the images are CT scans rather
than natural images.

\subsection{Data Loading Pipeline}

The \texttt{LDCTDataset} PyTorch \texttt{Dataset} class implements the complete image
loading and pre-processing pipeline for a given split. During \texttt{\_\_getitem\_\_},
the following operations are performed in sequence:

\begin{enumerate}
    \item \textbf{Image loading}: The 32-bit float TIFF file is opened with PIL,
          converted to a NumPy array of shape $(512, 512)$, and clipped to $[0, 1]$.
    \item \textbf{Augmentation} (training split only): stochastic spatial and
          photometric transformations are applied (see Section~\ref{sec:augmentation}).
    \item \textbf{Multi-scale pyramid construction}: The image is resized to each
          target scale using bicubic interpolation.
    \item \textbf{Patch extraction}: Each resized image is divided into
          non-overlapping $32 \times 32$ patches, and the spatial coordinates
          (scale index, row index, column index) of each patch are recorded.
    \item \textbf{Channel replication}: Each greyscale patch is replicated into
          three identical channels to produce an RGB-compatible tensor.
    \item \textbf{Tensor conversion}: All patches and coordinates are converted to
          PyTorch tensors.
\end{enumerate}

The \texttt{DataLoader} is configured with \texttt{batch\_size=10},
\texttt{shuffle=True} for training and \texttt{shuffle=False} for validation and test,
\texttt{num\_workers=0} for Windows-compatible multiprocessing, \texttt{pin\_memory=True}
for efficient GPU transfer, and \texttt{drop\_last=True} for training.


% ============================================================
\section{Proposed Model Architecture: CT-MUSIQ}
\label{sec:architecture}

\subsection{Architecture Overview}

CT-MUSIQ is a Transformer-based no-reference image quality assessor that operates on
multi-scale non-overlapping image patches. Inspired by the MUSIQ model of Ke et al.\
\cite{ke2021musiq}, which demonstrated that a Vision Transformer processing patches
from images at multiple resolutions simultaneously can learn richer quality-relevant
representations than single-scale CNNs, CT-MUSIQ adapts this paradigm for the specific
challenges of brain CT quality assessment: single-channel greyscale images, small
dataset size (700 training samples), and tight GPU memory constraints.

The architecture consists of five principal components, executed in sequence during each
forward pass:

\begin{enumerate}
    \item \textbf{Multi-scale patch tokenisation}: The input CT image is resized to
          two scales ($224 \times 224$ and $384 \times 384$) and divided into
          non-overlapping $32 \times 32$ patches, yielding 49 and 144 patches
          respectively, for a total of 193 patch tokens per image.
    \item \textbf{Patch Embedding}: Each patch is projected from pixel space
          ($3 \times 32 \times 32$) to a 768-dimensional embedding space using a shared
          2D convolutional projection layer initialised with pretrained ViT-B/32 weights.
    \item \textbf{Hash-based Positional Encoding}: Spatial position and scale identity
          are injected into each token via a sum of three learned embedding lookups.
    \item \textbf{Transformer Encoder}: A 12-layer, 8-head Transformer encoder
          processes the full sequence of 193 patch tokens prepended with a learnable
          \texttt{[CLS]} token (194 tokens total), performing global self-attention.
    \item \textbf{Prediction Heads}: A two-layer MLP global head reads the
          \texttt{[CLS]} token to produce the final quality score. Two auxiliary
          per-scale heads read scale-averaged token representations to produce
          scale-specific scores used only during training for the KL consistency loss.
\end{enumerate}

\subsection{Multi-Scale Patch Tokenisation}

The motivation for multi-scale representation in IQA is well established: quality
degradations such as noise, blur, and compression artefacts manifest differently at
different spatial frequencies and resolutions. A patch at a coarse scale captures global
structural context, while a patch at a fine scale captures local texture and noise
statistics. Processing both simultaneously enables the model to integrate quality cues
across spatial scales.

Two target scales are used: $S_1 = 224$~pixels and $S_2 = 384$~pixels. Each CT image
($512 \times 512$~pixels) is independently resized to each target scale using bicubic
interpolation. Each resized image is then divided into a regular grid of non-overlapping
$32 \times 32$ pixel patches:

\begin{itemize}
    \item At scale 224: grid size $= 224 \div 32 = 7$, yielding $7 \times 7 = 49$ patches.
    \item At scale 384: grid size $= 384 \div 32 = 12$, yielding $12 \times 12 = 144$ patches.
    \item Total patch count: $49 + 144 = 193$ patches per image.
\end{itemize}

The choice of patch size $P = 32$ is motivated by compatibility with the pretrained
Vision Transformer ViT-B/32, which uses $32 \times 32$ patches. Patches from both
scales are concatenated into a single sequence (scale-0 patches first, then scale-1
patches).

\subsection{Spatial Coordinate Representation}

Each patch is assigned a three-dimensional coordinate vector $(s_i, r_i, c_i)$ where
$s_i \in \{0, 1\}$ is the scale index, $r_i \in \{0, \ldots, g_s - 1\}$ is the row
index within the grid of scale $s$, and $c_i \in \{0, \ldots, g_s - 1\}$ is the column
index, with $g_s = S_s \div P$ being the grid size at scale $s$. For the 193 total
patches, the coordinate array has shape $[193, 3]$. These coordinates are stored as
integer tensors and consumed by the hash-based positional encoding module.

\subsection{Grayscale-to-RGB Channel Replication}

The pretrained ViT-B/32 model was trained on ImageNet RGB images, so its patch embedding
convolution has shape \texttt{[768, 3, 32, 32]}. Brain CT images are inherently
single-channel (greyscale). To exploit the pretrained patch embedding weights without
modification, each greyscale patch of shape $[1, 32, 32]$ is replicated into three
identical channels to produce a tensor of shape $[3, 32, 32]$:

\begin{equation}
    \mathbf{p}_{\text{RGB}} = [p, p, p]_{\text{channels}}
\end{equation}

This replication means the embedding layer can be initialised with meaningful pretrained
features rather than requiring random initialisation, which is especially valuable given
the small training dataset of 700 images. No artificial colour information is introduced.

\subsection{Patch Embedding Module}

The \texttt{PatchEmbedding} module is a single 2D convolutional layer with kernel size
equal to the patch size ($P = 32$) and stride equal to $P$. For a batch of patches of
shape $[B \cdot N, 3, 32, 32]$, the convolution produces $[B \cdot N, 768, 1, 1]$.
After flattening, each patch is represented as a 768-dimensional embedding vector.
Formally, for patch tensor $\mathbf{p} \in \mathbb{R}^{3 \times 32 \times 32}$:

\begin{equation}
    \mathbf{e} = \text{Flatten}\bigl(\text{Conv2d}(\mathbf{p})\bigr) \in \mathbb{R}^{768}
\end{equation}

The projection weights are initialised with pretrained ViT-B/32 weights from the
\textit{timm} library (\texttt{vit\_base\_patch32\_224}).

\subsection{Hash-Based Positional Encoding}

Standard sinusoidal or learned positional encodings in Vision Transformers assume a
fixed sequence length. CT-MUSIQ processes variable-composition sequences of 193 tokens
(49 from scale~0 plus 144 from scale~1), and the distinction between scales is crucial
for the model to exploit scale-specific features.

The custom \texttt{HashPositionalEncoding} module addresses this by decomposing position
into three orthogonal components, each represented by an independent learnable embedding
table:

\begin{equation}
    \mathbf{pos}_i = \mathbf{E}_{\text{scale}}(s_i) + \mathbf{E}_{\text{row}}(r_i)
                   + \mathbf{E}_{\text{col}}(c_i)
\end{equation}

where $\mathbf{E}_{\text{scale}} \in \mathbb{R}^{n_s \times d}$,
$\mathbf{E}_{\text{row}} \in \mathbb{R}^{G_{\max} \times d}$, and
$\mathbf{E}_{\text{col}} \in \mathbb{R}^{G_{\max} \times d}$ are three learned embedding
matrices; $n_s = 2$ is the number of scales; $G_{\max} = 13$ is the maximum grid
dimension; and $d = 768$ is the embedding dimension. All embedding tables are initialised
with zero-mean Gaussian noise with standard deviation 0.02.

This decomposition has several desirable properties. First, it handles the
variable-length multi-scale sequence without requiring padding. Second, the scale
embedding $\mathbf{E}_{\text{scale}}(s)$ allows the model to learn scale-specific
biases. Third, the factored structure dramatically reduces the embedding table size to
$(2 + 13 + 13) \times 768 = 21{,}504$ parameters versus $2 \times 13 \times 13 \times
768 = 328{,}704$ for a full 3D lookup table. The positional encoding is added
element-wise to the patch embeddings:

\begin{equation}
    \mathbf{z}_i = \mathbf{e}_i + \mathbf{pos}_i, \quad i = 1, \ldots, N
\end{equation}

\subsection{Learnable \texttt{[CLS]} Token}

Following the standard Vision Transformer design, a learnable global classification
token (\texttt{[CLS]} token) is prepended to the patch token sequence. The \texttt{[CLS]}
token is a single parameter vector $\boldsymbol{\theta}_{\text{cls}} \in \mathbb{R}^{768}$,
initialised from $\mathcal{N}(0, 0.02^2)$ and replicated across the batch dimension
during the forward pass. After the Transformer encoder processes the full sequence of
$N + 1 = 194$ tokens, the \texttt{[CLS]} token at position 0 in the output sequence
aggregates global context from all 193 patch tokens via the self-attention mechanism.

\subsection{Transformer Encoder}

The Transformer encoder is a stack of $L = 12$ \texttt{TransformerEncoderLayer} modules.
Each layer operates on the full sequence of 194 tokens and performs the following
operations:

\begin{enumerate}
    \item \textbf{Pre-Layer Normalisation}: Layer normalisation is applied before each
          sub-layer (``norm-first'' configuration), which improves training stability
          when fine-tuning on small datasets.
    \item \textbf{Multi-Head Self-Attention}: With $H = 8$ attention heads and
          $d_{\text{model}} = 768$, each head has per-head dimension
          $d_k = 768 / 8 = 96$. Global self-attention is performed over all 194 tokens,
          enabling the model to capture long-range spatial dependencies across both
          scales simultaneously.
    \item \textbf{Feed-Forward Network}: A position-wise FFN with hidden dimension
          $d_{\text{ff}} = 3072 = 4 \times d_{\text{model}}$, using GELU activation and
          dropout $p = 0.05$.
    \item \textbf{Residual Connections}: Standard residual connections around each
          sub-layer.
\end{enumerate}

The dropout rate is deliberately reduced to $p = 0.05$ (versus the standard 0.1 in ViT)
to mitigate overfitting on the small 700-sample training set. The encoder weights are
initialised from pretrained ViT-B/32 weights, covering all attention projection weights,
layer normalisation parameters, and FFN weights for all 12 layers.

\subsection{Prediction Heads}

\textbf{Global Head:} The \texttt{[CLS]} token output
$\mathbf{c} \in \mathbb{R}^{768}$ is passed through a two-layer MLP:

\begin{equation}
    \hat{q} = W_2 \cdot \text{Dropout}\bigl(\text{GELU}(W_1 \mathbf{c})\bigr)
\end{equation}

where $W_1 \in \mathbb{R}^{384 \times 768}$ and $W_2 \in \mathbb{R}^{1 \times 384}$.
The output is a single scalar $\hat{q} \in \mathbb{R}$, which is the predicted quality
score at inference time.

\textbf{Per-Scale Heads:} For each scale $s \in \{0, 1\}$, the patch token outputs
corresponding to that scale are extracted and mean-pooled along the patch dimension to
produce a scale-specific representation $\mathbf{c}_s \in \mathbb{R}^{768}$. This is
passed through an identical two-layer MLP to produce a per-scale quality score
$\hat{q}_s$. The per-scale scores are used exclusively during training to compute the
scale-consistency KL divergence loss (Section~\ref{sec:loss}).


% ============================================================
\section{Loss Function}
\label{sec:loss}

CT-MUSIQ is trained with a composite loss function combining three terms: a primary
regression loss, a scale-consistency regularisation loss, and a pairwise ranking loss.
The combined objective is:

\begin{equation}
    \mathcal{L}_{\text{total}} = \mathcal{L}_{\text{MSE}}
                               + \lambda_{\text{KL}} \cdot \mathcal{L}_{\text{KL}}
                               + \lambda_{\text{rank}} \cdot \mathcal{L}_{\text{rank}}
    \label{eq:total_loss}
\end{equation}

where $\lambda_{\text{KL}} = 0.10$ (determined by ablation study,
Section~\ref{sec:ablation}) and $\lambda_{\text{rank}} = 0.5$.

\subsection{Mean Squared Error Loss}

The primary regression objective is the standard mean squared error (MSE) between the
global predicted score $\hat{q}$ and the ground truth radiologist score $q$:

\begin{equation}
    \mathcal{L}_{\text{MSE}} = \frac{1}{B} \sum_{i=1}^{B} \bigl(\hat{q}_i - q_i\bigr)^2
\end{equation}

MSE penalises large errors quadratically, encouraging the model to be accurate for
images whose quality is clearly very high or very low, which is clinically most important.

\subsection{Scale-Consistency KL Divergence Loss}

The key regularisation innovation in CT-MUSIQ is the \textbf{scale-consistency KL
divergence loss}, which penalises disagreement between the per-scale quality predictions
($\hat{q}_s$) and the global quality prediction ($\hat{q}$). The intuition is that if
both scales observe the same image at different resolutions, their quality assessments
should be broadly consistent. A model that learns scale-specific shortcuts is learning
artefacts rather than true quality features.

Computing KL divergence between scalar scores requires converting them to probability
distributions. This is accomplished via a \textbf{Gaussian kernel smoothing} procedure
over $K = 20$ discrete bins uniformly spanning the score range $[0, 4]$. For a scalar
score $\hat{q}$ and bin centres $\{b_k\}_{k=1}^{K}$:

\begin{align}
    \tilde{p}_k(\hat{q}) &= \exp\!\left(-\frac{(\hat{q} - b_k)^2}{2\sigma^2}\right) \\
    p_k(\hat{q})          &= \frac{\tilde{p}_k(\hat{q})}{\displaystyle\sum_{j=1}^{K}
                             \tilde{p}_j(\hat{q}) + \varepsilon}
\end{align}

with $\sigma = 0.5$ and $\varepsilon = 10^{-8}$ for numerical stability. The result is a
soft probability distribution $\mathbf{p}(\hat{q}) \in \Delta^{K-1}$ that peaks at the
bin nearest to $\hat{q}$ and decays smoothly away. The scale-consistency loss is then:

\begin{equation}
    \mathcal{L}_{\text{KL}} = \frac{1}{n_s} \sum_{s=0}^{n_s - 1}
    D_{\text{KL}}\!\Bigl(\mathbf{p}(\hat{q}_s) \;\|\; \mathbf{p}(\hat{q})\Bigr)
\end{equation}

where $D_{\text{KL}}$ denotes the KL divergence computed using log-probabilities of the
scale distribution against the probability of the global distribution, with batch-mean
reduction. The global prediction is treated as the target distribution.

\subsection{Pairwise Margin Ranking Loss}

The MSE loss optimises absolute score accuracy but does not explicitly penalise
incorrect pairwise orderings. Since the LDCTIQAC 2023 evaluation metrics (PLCC, SROCC,
KROCC) all measure rank-based correlation, training with an explicit ranking objective
is expected to improve performance. The \textbf{pairwise margin ranking loss} operates
on all $B(B-1)$ ordered pairs $(i, j)$ in each batch of size $B$:

\begin{equation}
    \mathcal{L}_{\text{rank}} = \frac{1}{B^2} \sum_{i=1}^{B} \sum_{j=1}^{B}
    \max\!\Bigl(0,\; m - \operatorname{sgn}(q_i - q_j) \cdot (\hat{q}_i - \hat{q}_j)\Bigr)
\end{equation}

where $m = 0.1$ is the margin. The term $\operatorname{sgn}(q_i - q_j)$ is $+1$ if
image~$i$ truly has higher quality than image~$j$, $-1$ if lower, and 0 if equal. The
loss encourages the predicted score difference $\hat{q}_i - \hat{q}_j$ to have the same
sign as the true difference and to exceed the margin $m$, directly optimising relative
ordering that translates to improved rank correlation metrics.


% ============================================================
\section{Training Procedure}
\label{sec:training}

\subsection{Two-Stage Training Schedule}

CT-MUSIQ employs a two-stage training strategy designed to stabilise fine-tuning of
pretrained ViT weights on a small medical imaging dataset.

\textbf{Stage~1 (Epochs 1--5): Frozen Encoder Warm-Up.}
In the first stage, the Transformer encoder weights and the \texttt{[CLS]} token
parameter are frozen. Only the patch embedding layer, hash positional encoding, and
prediction heads are trained. The learning rate is set to
$\text{LR}_{\text{stage1}} = 3 \times 10^{-4}$, which is relatively high to allow the
randomly-initialised CT-specific components to reach a useful parameter regime quickly.
The rationale is to prevent the noisy gradients from randomly-initialised prediction
heads from corrupting the carefully pretrained attention patterns during the earliest
training steps.

\textbf{Stage~2 (Epochs 6--100): Full Fine-Tuning.}
Beginning at epoch~6, all parameters are unfrozen. The learning rate transitions through
two sub-phases within Stage~2:

\begin{enumerate}
    \item \textbf{Linear Warm-Up (Epochs 6--8, \texttt{STAGE2\_WARMUP\_EPOCHS} = 3)}:
          The learning rate is linearly increased from
          $\text{LR}_{\text{stage2}} \cdot (1/3)$ to
          $\text{LR}_{\text{stage2}} = 5 \times 10^{-5}$ over 3 epochs. This warm-up
          prevents the abrupt large-gradient update that would occur if the encoder
          were suddenly fully unfrozen at full learning rate.
    \item \textbf{Cosine Annealing (Epochs 9--100)}: After the warm-up, a cosine
          annealing scheduler gradually decays the learning rate from
          $5 \times 10^{-5}$ to $\text{LR}_{\min} = 10^{-6}$:
          \begin{equation}
              \text{LR}(t) = \text{LR}_{\min} + \frac{1}{2}
              (\text{LR}_{\max} - \text{LR}_{\min})
              \left(1 + \cos\!\left(\frac{\pi t}{T_{\max}}\right)\right)
          \end{equation}
          where $t$ is the number of epochs elapsed since the start of the cosine phase
          and $T_{\max}$ is the total number of cosine epochs.
\end{enumerate}

\subsection{Optimiser}

The AdamW optimiser (Adam with decoupled weight decay) is used throughout training with
weight decay coefficient $\lambda_{\text{wd}} = 0.01$. AdamW is preferred over standard
Adam for fine-tuning pretrained Transformers because decoupled weight decay provides
more principled regularisation: it applies decay uniformly to all parameters rather than
scaling it by the adaptive learning rate per parameter.

\subsection{Mixed Precision Training}

To fit the 12-layer ViT-B-based model with batch size~10 within 6~GB of GPU VRAM,
\textbf{automatic mixed precision (AMP)} training is employed using PyTorch's
\texttt{torch.cuda.amp} framework. AMP operates in two modes: the forward pass and loss
computation use half-precision (fp16) tensors for activations, while the backward pass
and parameter updates are performed in full precision (fp32). The approximate VRAM
breakdown with AMP is as follows:
model weights~$\approx$~1.5~GB, activations~$\approx$~1.5--2~GB, and optimiser
states~$\approx$~1.5~GB, for a total well within the 6~GB budget.
Mixed precision saves approximately 40\% of VRAM relative to pure fp32 and provides
approximately 30\% speedup on the RTX~3060 GPU.

A \texttt{GradScaler} is used to dynamically scale the loss before the backward pass to
prevent fp16 gradient underflow. The scaler reduces the scale factor if overflow is
detected and gradually increases it otherwise.

\subsection{Gradient Clipping}

Gradient clipping is applied after the \texttt{GradScaler.unscale\_()} step, which
restores gradients to fp32 scale before clipping. The maximum gradient norm is
$\|\nabla\|_{\max} = 1.0$, applied to all trainable parameters jointly using the global
$\ell_2$ norm. This prevents the occasional large gradient updates that can occur at the
Stage~1 to Stage~2 transition.

\subsection{Exponential Moving Average}

An \textbf{Exponential Moving Average (EMA)} of the model weights is maintained during
training with decay factor $\rho = 0.999$:

\begin{equation}
    \boldsymbol{\theta}_{\text{EMA}}^{(t)} = \rho \cdot \boldsymbol{\theta}_{\text{EMA}}^{(t-1)}
    + (1 - \rho) \cdot \boldsymbol{\theta}^{(t)}
\end{equation}

During validation at the end of each epoch, the EMA weights are temporarily substituted
into the model (the live weights are saved to a backup), the validation pass is
performed, and then the live weights are restored. The checkpoint saved when validation
performance improves stores the EMA weights. This ensures the saved model benefits from
the smoothing effect of EMA, which tends to have better generalisation properties.

\subsection{Early Stopping}

Training is monitored against the validation aggregate score
$|\text{PLCC}| + |\text{SROCC}| + |\text{KROCC}|$. If the validation aggregate does
not improve for \texttt{PATIENCE} $= 20$ consecutive epochs, training is halted early
to prevent overfitting. In all conducted full training runs, early stopping was not
triggered before epoch~100, indicating that the model continued improving through the
full training budget.

\subsection{Reproducibility}

All random number generators are seeded with \texttt{SEED} $= 42$ at the beginning of
each training run via \texttt{torch.manual\_seed()},
\texttt{torch.cuda.manual\_seed\_all()}, \texttt{numpy.random.seed()}, and
\texttt{torch.backends.cudnn.deterministic = True}.


% ============================================================
\section{Data Augmentation}
\label{sec:augmentation}

Given the small training set of 700 images, data augmentation is essential to improve
generalisation and prevent overfitting. CT imaging places unique constraints on
acceptable augmentations: unlike natural images, CT images are calibrated physical
measurements in Hounsfield Units, and their noise texture and structural patterns are
diagnostically meaningful. Aggressive colour distortions or large geometric deformations
would alter the perceived quality in ways that do not correspond to any real imaging
artefact, creating label noise. Therefore, augmentations are carefully chosen to be mild
and clinically plausible. The following augmentations are applied stochastically during
training only (never during validation or testing):

\begin{itemize}
    \item \textbf{Horizontal Flip} ($p = 0.5$): The CT image is flipped left-to-right.
          Axial brain CT slices are broadly symmetric across the mid-sagittal plane,
          making horizontal flip a valid augmentation.
    \item \textbf{Vertical Flip} ($p = 0.5$): The CT image is flipped top-to-bottom.
          Qualitative characteristics of noise and artefacts are isotropic in axial
          slices, making vertical flip reasonable.
    \item \textbf{Small Random Rotation} ($p = 0.5$, angle $\in [-7°, +7°]$): Applied
          using bilinear interpolation in float-mode PIL to avoid 8-bit quantisation
          artefacts.
    \item \textbf{Random Crop and Resize} ($p = 0.5$, crop fraction $\in [0.85, 1.0]$):
          A random crop retaining 85--100\% of the original image area is taken and
          bicubically resized back to $512 \times 512$. This simulates slight variation
          in the field of view.
    \item \textbf{Mild Brightness and Contrast Jitter} (delta $\in [-0.05, +0.05]$):
          Brightness is perturbed by a random additive offset and contrast by a
          multiplicative factor, both sampled from $[-0.05, +0.05]$.
    \item \textbf{Mild Additive Gaussian Noise} (std $\in [0.005, 0.02]$): Random
          Gaussian noise with standard deviation uniformly sampled from
          $[0.005, 0.02]$ is added to simulate scanner noise variability.
\end{itemize}

All augmentations are applied in float32 precision, and the result is clipped to
$[0, 1]$ after all transformations.


% ============================================================
\section{Evaluation Protocol}
\label{sec:evaluation}

\subsection{Evaluation Metrics}

The evaluation metrics precisely follow the LDCTIQAC 2023 challenge protocol as
described in Lee et al.\ \cite{lee2025ldctiqac}. Three correlation coefficients are
computed between the predicted scores and the ground truth scores:

\textbf{Pearson Linear Correlation Coefficient (PLCC):}
\begin{equation}
    \text{PLCC} = \frac{\sum_{i}(\hat{q}_i - \bar{\hat{q}})(q_i - \bar{q})}
                       {\sqrt{\sum_{i}(\hat{q}_i - \bar{\hat{q}})^2
                              \sum_{i}(q_i - \bar{q})^2}}
\end{equation}

\textbf{Spearman Rank-Order Correlation Coefficient (SROCC):}
\begin{equation}
    \text{SROCC} = 1 - \frac{6\sum_{i} d_i^2}{N_{\text{test}}(N_{\text{test}}^2 - 1)}
\end{equation}
where $d_i$ is the difference in ranks between $\hat{q}_i$ and $q_i$.

\textbf{Kendall Rank-Order Correlation Coefficient (KROCC):}
\begin{equation}
    \text{KROCC} = \frac{n_c - n_d}{\frac{1}{2}N_{\text{test}}(N_{\text{test}} - 1)}
\end{equation}
where $n_c$ is the number of concordant pairs and $n_d$ is the number of discordant pairs.

\textbf{Aggregate Score:}
\begin{equation}
    \text{Aggregate} = |\text{PLCC}| + |\text{SROCC}| + |\text{KROCC}|
    \label{eq:aggregate}
\end{equation}

The aggregate score is the primary ranking metric used in the LDCTIQAC 2023 challenge
leaderboard, with maximum possible value of 3.0 (perfect correlation on all three
metrics).

\subsection{Test-Time Augmentation}

Test-time augmentation (TTA) with horizontal flip averaging is implemented as an optional
inference strategy. For each test image, the model is run twice: once on the original
patch sequence and once on a horizontally-flipped version. The horizontal flip remaps
column coordinates as:

\begin{equation}
    c_{\text{flipped}} = (g_s - 1) - c_{\text{original}}
\end{equation}

where $g_s$ is the grid size at scale $s$. The final prediction is the average of the
two forward pass outputs, exploiting the approximate left-right symmetry of brain CT
images.

\subsection{Isotonic Regression Calibration}

An optional post-processing calibration step using \textbf{isotonic regression} is
provided. Isotonic regression fits a non-decreasing step function mapping raw predictions
to calibrated outputs, fitted on the validation set:

\begin{equation}
    \hat{q}^*_{\text{test}} = \text{ISO}(\hat{q}_{\text{test}})
\end{equation}

where \texttt{ISO} is the isotonic regression function fitted on validation set pairs
$\{(\hat{q}_{\text{val},i}, q_{\text{val},i})\}$. Because isotonic regression is a
monotone transformation, it preserves rank order exactly (SROCC and KROCC are unchanged)
while potentially improving PLCC by correcting systematic scale or bias errors.


% ============================================================
\section{Comparison Baseline: Agaldran Combo}
\label{sec:baseline}

For a fair within-paper comparison, a strong baseline model is trained and evaluated
under identical conditions. The baseline, termed \textbf{Agaldran Combo}, is motivated
by the top-performing method in the LDCTIQAC 2023 challenge (the ``agaldran'' team,
ranked \#1 with Aggregate $= 2.7427$ on the published leaderboard), which combined a
Swin Transformer with a ResNet backbone~\cite{lee2025ldctiqac}.

The baseline model is a \textbf{score-level ensemble} of two independently pretrained
backbones:

\begin{enumerate}
    \item \textbf{Swin Transformer Tiny (Swin-T)}: Loaded from \textit{timm}
          (\texttt{swin\_tiny\_patch4\_window7\_224}) with ImageNet-1K pretrained
          weights. The classification head is removed and replaced with a quality
          prediction head: \texttt{Linear(768, 256)} $\to$ \texttt{ReLU} $\to$
          \texttt{Dropout(0.05)} $\to$ \texttt{Linear(256, 1)}.
    \item \textbf{ResNet-50}: Loaded from \textit{torchvision} with ImageNet-1K
          pretrained weights. The fully connected classification layer is replaced by a
          quality prediction head: \texttt{Linear(2048, 256)} $\to$ \texttt{ReLU} $\to$
          \texttt{Dropout(0.05)} $\to$ \texttt{Linear(256, 1)}.
\end{enumerate}

The final prediction is the arithmetic mean of the two backbones' quality scores:

\begin{equation}
    \hat{q}_{\text{combo}} = \frac{\hat{q}_{\text{Swin-T}} + \hat{q}_{\text{ResNet50}}}{2}
\end{equation}

The baseline accepts full CT images as input (not patches) of shape $[B, 3, H, W]$.
A \texttt{build\_baseline\_images\_from\_patches} utility function reconstructs the
original scale-0 ($224 \times 224$) image from the scale-0 patch tokens by placing each
patch back into its grid position according to the coordinate tensor.

The baseline is trained under the same loss function (\texttt{CTMUSIQLoss}), batch
size~(10), and dataset split as CT-MUSIQ. Unlike CT-MUSIQ, the baseline does not employ
the two-stage training schedule; it is trained end-to-end from the start with a constant
learning rate of $3 \times 10^{-4}$ and the same early stopping patience of 20 epochs.


% ============================================================
\section{Ablation Study Design}
\label{sec:ablation}

To validate the individual contributions of CT-MUSIQ's proposed components, five
ablation experiments are designed. Each ablation modifies a single aspect of the full
model while keeping all other settings constant. All ablation runs are trained for
30~epochs.

\begin{table}[h]
    \centering
    \caption{Ablation experiment configurations.}
    \label{tab:ablations}
    \begin{tabular}{llcc}
        \toprule
        Config & Description & Scales & $\lambda_{\text{KL}}$ \\
        \midrule
        A1 & Single-scale baseline (224 only, no KL) & [224]         & 0.0  \\
        A2 & Two scales, no KL loss                  & [224, 384]    & 0.0  \\
        A3 & Two scales, KL $\lambda = 0.05$ (weak)  & [224, 384]    & 0.05 \\
        A4 & Two scales, KL $\lambda = 0.10$ (optimal) & [224, 384]  & 0.10 \\
        A5 & Two scales, KL $\lambda = 0.20$ (strong) & [224, 384]   & 0.20 \\
        A6 & Three scales, KL $\lambda = 0.10$        & [224, 384, 512] & 0.10 \\
        \bottomrule
    \end{tabular}
\end{table}

\noindent
The ablation comparisons address the following questions:
\begin{itemize}
    \item \textbf{A1 vs.\ A2}: Is multi-scale representation beneficial? If A2
          outperforms A1, the second scale (384) provides useful complementary features.
    \item \textbf{A2 vs.\ A3/A4/A5}: Does the scale-consistency KL loss improve
          performance, and what is the optimal weight? The sweep across
          $\lambda_{\text{KL}} \in \{0.05, 0.10, 0.20\}$ identifies the best balance
          between MSE accuracy and cross-scale consistency.
    \item \textbf{A4 vs.\ A6}: Does a three-scale pyramid (adding 512~px,
          yielding $16 \times 16 = 256$ additional patches and 449 total tokens) provide
          further improvement, given sufficient VRAM?
\end{itemize}


% ============================================================
\section{Implementation Details}
\label{sec:implementation}

\subsection{Software Environment}

The project is implemented in Python~3.10 using PyTorch~2.x as the deep learning
framework. Key dependencies include:
\textit{timm}~($\geq$0.9.0) for pretrained ViT-B/32 and Swin-T weights,
\textit{torchvision}~($\geq$0.15) for ResNet-50 weights,
\textit{scipy} for correlation metric computation,
\textit{scikit-learn} for isotonic regression calibration,
\textit{PIL/Pillow} for float-mode image manipulation,
\textit{numpy} for array operations during pre-processing and augmentation, and
\textit{pandas} for training log CSV management.

All hyperparameters are centralised in \texttt{config.py}, which acts as the single
source of truth for the entire project. No hyperparameter is hard-coded in training,
evaluation, or dataset scripts.

\subsection{Hardware}

All experiments are conducted on a workstation equipped with an \textbf{NVIDIA RTX~3060
Laptop GPU} with 6~GB GDDR6 VRAM, a mobile Intel Core~i7 CPU, and 16~GB of system RAM.
Training one epoch on the 700-image training set takes approximately \textbf{31--40
seconds} at convergence. A complete 100-epoch training run takes approximately
\textbf{50--55 minutes}.

\subsection{Model Parameter Count}

The CT-MUSIQ model has approximately \textbf{88.3 million parameters} in total:
the ViT-B encoder ($\sim$85M), patch embedding ($\sim$2.3M), hash positional encoding
($\sim$21.5K), global and per-scale prediction heads ($\sim$887K), and the
\texttt{[CLS]} token (768 parameters). At fp32 precision, the model occupies
approximately 353~MB on disk. The full checkpoint including EMA shadow weights and
optimiser states is approximately 1~GB.

\subsection{Training Convergence}

Training log data from the final CT-MUSIQ run (100 epochs, full two-stage schedule)
reveals smooth convergence behaviour. In Stage~1 (epochs 1--5), the validation aggregate
rises from approximately 0 to $\sim$1.26 as the randomly-initialised prediction heads
learn a crude quality mapping. At the Stage~2 transition (epoch~6), the aggregate rises
steeply: $\sim$1.88 by epoch~6, $\sim$2.10 by epoch~10, above 2.50 by epoch~30, and
converging to a plateau near 2.76 in the final epochs. The training MSE decays from
approximately 1.43 at epoch~1 to approximately 0.12 by epoch~100.

\subsection{Final Experimental Results}

The final model checkpoint achieves the following performance on the \textbf{held-out
test set} (images 0900--0999, $N_{\text{test}} = 100$):

\begin{table}[h]
    \centering
    \caption{Test set results comparison with published LDCTIQAC 2023 leaderboard
             results~\cite{lee2025ldctiqac}.}
    \label{tab:results}
    \begin{tabular}{lcccc}
        \toprule
        Method & Aggregate & PLCC & SROCC & KROCC \\
        \midrule
        agaldran~\cite{lee2025ldctiqac}   & 2.7427 & 0.9491 & 0.9495 & 0.8440 \\
        RPI\_AXIS~\cite{lee2025ldctiqac}  & 2.6843 & 0.9434 & 0.9414 & 0.7995 \\
        CHILL@UK~\cite{lee2025ldctiqac}   & 2.6719 & 0.9402 & 0.9387 & 0.7930 \\
        FeatureNet~\cite{lee2025ldctiqac} & 2.6550 & 0.9362 & 0.9338 & 0.7851 \\
        Team Epoch~\cite{lee2025ldctiqac} & 2.6202 & 0.9278 & 0.9232 & 0.7691 \\
        gabybaldeon~\cite{lee2025ldctiqac} & 2.5671 & 0.9143 & 0.9096 & 0.7432 \\
        SNR baseline~\cite{lee2025ldctiqac} & 2.4026 & 0.8226 & 0.8748 & 0.7052 \\
        BRISQUE~\cite{lee2025ldctiqac}    & 2.1219 & 0.7500 & 0.7863 & 0.5856 \\
        \midrule
        \textbf{CT-MUSIQ (ours)} & \textbf{2.7682} & \textbf{0.9607} & \textbf{0.9605} & \textbf{0.8469} \\
        \bottomrule
    \end{tabular}
\end{table}

CT-MUSIQ surpasses the top published result in the LDCTIQAC 2023 challenge leaderboard
on all three correlation metrics and achieves the highest aggregate score, demonstrating
state-of-the-art perceptual quality assessment for brain low-dose CT images.


% ============================================================
\section{Summary}
\label{sec:methodology_summary}

This chapter has described the complete methodology of CT-MUSIQ, a no-reference image
quality assessment system for brain low-dose CT imaging. The methodology can be
summarised across its six principal contributions:

\begin{enumerate}
    \item \textbf{Clinically-motivated CT windowing}: Applying the brain soft-tissue
          window (level $= 40$, width $= 80$~HU) instead of the challenge's original
          abdominal window, aligning pre-processing with radiological practice for brain
          CT and maximising the discriminability of quality-relevant tissue contrast.
    \item \textbf{Multi-scale patch tokenisation with hash positional encoding}:
          Resizing each image to 224 and 384~pixels, extracting $32 \times 32$ patches
          to produce 193 total tokens, and injecting spatial and scale identity via a
          novel three-component additive embedding scheme that handles variable-length
          multi-scale sequences without the limitations of standard fixed-length
          positional encodings.
    \item \textbf{Pretrained ViT-B/32 transfer with greyscale-to-RGB replication}:
          Leveraging ImageNet-pretrained Vision Transformer weights (both patch embedding
          and full 12-layer encoder) for CT quality assessment, adapting the
          single-channel greyscale input by channel replication to achieve compatibility
          with the pretrained RGB weights.
    \item \textbf{Composite loss function}: Combining MSE regression loss,
          KL divergence-based scale-consistency regularisation with Gaussian-kernel
          score-to-distribution conversion, and pairwise margin ranking loss into a
          unified training objective that simultaneously optimises absolute accuracy,
          cross-scale agreement, and relative ordering correctness.
    \item \textbf{Two-stage training with EMA and mixed precision}: A structured
          Stage~1 (frozen encoder) $\to$ Stage~2 (full fine-tuning with warmup and
          cosine annealing) schedule, supported by exponential moving average smoothing
          of weights, gradient clipping, and mixed-precision training for memory
          efficiency on constrained GPU hardware.
    \item \textbf{Rigorous evaluation and ablation}: Evaluation on the LDCTIQAC 2023
          test set using the official challenge metrics (PLCC, SROCC, KROCC, Aggregate),
          comparison with published leaderboard results, a trained baseline (Agaldran
          Combo), and a systematic ablation study quantifying the independent
          contributions of multi-scale representation and scale-consistency
          regularisation.
\end{enumerate}
